\section{Post-review Development: Invariant Moments and a Bounded Viewing Law}
\label{app:momentdevelopment}
This section records subsequent method development, including failed gates.
It is not a new blinded confirmation or a claim of experimental calibration.
Bispectra and moment inference are established prior art
\citep{perry2019mra,janson2026mps}; the results below do not establish a novel
uncertainty method. The original full-review conclusions remain unresolved.

\subsection{Translation-invariant statistics and exact noise variance}
For nonzero, distinct Fourier coordinates without opposite pairs and
$q_a+q_b=q_c$, use $(|Y_j|^2-2v)/2$ together with
$\operatorname{Re}(Y_aY_b\overline{Y_c})/2$ and its imaginary counterpart.
All three indices in a triad are distinct. A common continuous translation
cancels from each product. Independent real and imaginary Gaussian noise
coordinates of variance $v$ make their expectations the noise-free moments.
This does not hold automatically under windowing, unknown correlations or
noncircular noise. Normalization to unit marginal variance at pure unit noise
does not whiten their signal-dependent joint covariance.

Let $f$ be a fixed real linear contrast of these moments. Its cubic polynomial
has constant Laplacian. The finite Gaussian--Hermite expansion therefore gives
\begin{align*}
\operatorname{Var}\{f(m+\epsilon)\}
&=v\|\nabla f(m)\|^2+\tfrac12v^2\|\nabla^2 f(m)\|_F^2\\
&\quad+\tfrac16v^3\|\nabla^3 f\|_F^2.
\end{align*}
Derivatives from overlapping triads must be added before squaring. This is a
classical polynomial identity, not a new general theorem. Independent complex
Gaussian moment enumeration agrees with 48 saved derivative-based variance
calculations within $3.6\times10^{-15}$.

\begin{table*}[t]
\centering\footnotesize
\caption{Post-review Monte Carlo development. Projected rejection probabilities on the Haar true-map alternative against a removed-map candidate. Columns at $\kappa=1$ vary particle count; the last three use $n=10^5$. Brackets are pointwise 95\% intervals for the $\kappa=1,n=10^5$ projection. All values are rounded to three decimals; 0.000/1.000 need not be exact probabilities. The full CSV retains all 192 projections and their intervals, including correct-null controls. Increasing $\kappa$ changes the test bound, not the held-out viewing law.}
\label{tab:momentmc}
\begin{tabular}{lllrrrrrr}
\toprule
Stack & Statistic & Amplitude & $n=10^3$ & $n=10^4$ & $n=10^5$ [95\%] & $\kappa=1.01$ & $\kappa=1.1$ & $\kappa=2$ \\
\midrule
10028 & Power & 1 & 0.865 & 1.000 & 1.000 [1.000, 1.000] & 1.000 & 0.268 & 0.000 \\
10028 & Power & [.9,1.1] & 0.022 & 0.004 & 0.000 [0.000, 0.001] & 0.000 & 0.000 & 0.000 \\
10028 & Power+bispectrum & 1 & 0.877 & 1.000 & 1.000 [1.000, 1.000] & 1.000 & 0.503 & 0.000 \\
10028 & Power+bispectrum & [.9,1.1] & 0.029 & 0.009 & 0.000 [0.000, 0.014] & 0.000 & 0.000 & 0.000 \\
10049 & Power & 1 & 0.624 & 1.000 & 1.000 [1.000, 1.000] & 1.000 & 0.000 & 0.000 \\
10049 & Power & [.9,1.1] & 0.166 & 0.700 & 1.000 [1.000, 1.000] & 0.994 & 0.000 & 0.000 \\
10049 & Power+bispectrum & 1 & 0.669 & 1.000 & 1.000 [1.000, 1.000] & 1.000 & 0.000 & 0.000 \\
10049 & Power+bispectrum & [.9,1.1] & 0.192 & 0.826 & 1.000 [1.000, 1.000] & 1.000 & 0.000 & 0.000 \\
10076 & All four contrasts & Both & \multicolumn{6}{l}{Zero directions; constant scores, no simulation and no rejection.} \\
\bottomrule
\end{tabular}
\end{table*}


\subsection{Finite hulls and continuous-orientation counterexamples}
The gate uses 220 frequencies and 512 frequency-selected complex triads,
giving 1,244 real moments, and the first saved CTF/noise profile. Each oracle
alternative averages the true-map moments over 4,160 saved orientations.
Power-only and combined moments are compared at amplitude one and on the
assumed interval $[.9,1.1]$, for all three geometries. Twelve exact true-map
controls and twelve removal fits complete; every removal fit reaches the
2,500-iteration limit. Feasible mixtures and separating directions give
finite-hull distance brackets without assuming convergence.

A separate 10,000-view catalog and 160 continuous local searches then test all
eight nonzero directions. Four zero directions on 10076 are retained as
uninformative. Both ranged-amplitude 10028 directions have verified null poses
whose expected scores exceed the alternative expectation. The four 10049
directions retain sampled margins $.0174$--$.0508$, without a global bound.
All eight worst searched poses independently agree with physical-cell sums
within $3.8\times10^{-12}$. Neither searches nor unconverged hull fits establish
continuous separation or impossibility for the moment family.

One explicit global envelope is computationally impractical. If the contrast
has geodesic second derivative bounded by $H$, and a rotation grid covers
$\mathrm{SO}(3)$ with radius $r$, its continuously amplitude-profiled maximum
is at most the profiled grid maximum plus $Hr^2/2$. To see this, fix a
maximizing amplitude and rotation; the rotational derivative vanishes there.
Integrating the second derivative to a neighboring grid point proves the
bound. Direct absolute-density radial moments give finite $H$, including
physical cell extent, but the declared Euler-grid construction needs
$1.1\times10^{13}$--$2.0\times10^{16}$ points to allocate half the remaining
sampled margin to this remainder. No grid is run. This is a failure of that
envelope, not an optimal-complexity lower bound.

\subsection{A different model: simulation under a bounded viewing density}
The following construction changes the nuisance model. Let $Q$ be a specified
viewing law, here continuous Haar measure, and assume each independent
particle's viewing distribution $P_i$ satisfies $dP_i/dQ\le\kappa$, for a
known $\kappa\ge1$. Noise and CTF are known; amplitudes belong to a fixed
interval $\mathcal A$. No arbitrary-realized-pose guarantee follows.
The score $f$ and threshold $c$ are fixed independently of calibration and
tested images. For a fixed candidate simulator $m(R)$, draw
\[
Z=\mathbf1\{\max_{a\in\mathcal A} f(a m(R)+\epsilon)>c\},
\qquad R\sim Q.
\]
The amplitude maximum uses both endpoints and all real stationary points
of a scalar cubic. Let $U$ be a one-sided Clopper--Pearson upper probability
from $M$ independent draws of $Z$, with failure probability $\delta$.
Set $p_0=\min(1,\kappa U)$. Reject the candidate when the number of $n$ tested
images with score exceeding $c$ reaches the smallest integer $k$ with
$\Pr\{\operatorname{Bin}(n,p_0)\ge k\}\le\alpha-\delta$.

\paragraph{Conditional-model validity.}
In exact arithmetic the joint false-rejection probability over calibration
and test data is at most $\alpha$. The amplitude envelope dominates every
allowed amplitude's event. The viewing density ratio increases its probability
by at most $\kappa$. Except on the calibration failure event, each independent
test indicator thus has probability at most $p_0$; coupling with independent
uniform variables gives binomial stochastic domination. A union bound adds
$\delta$ to the conditional $\alpha-\delta$ error. This is a composition of
classical binomial inference and domination. Monte Carlo tests and nuisance
maximization have substantial prior art \citep{dufour2006mc,bergerboos1994}.
We do not claim either cited paper's algorithm as this construction.

\subsection{Independent simulations and limits of the comparison}
All twelve saved contrasts and their previously fixed alternative-mean
thresholds are retained. Each nonzero stack receives 131,072 fresh continuous
Haar calibration draws and 131,072 independent held-out draws. The physical-cell
operator evaluates both true and removed maps; within each stage methods share
rotations and Gaussian noise. Both maps also serve as correctly specified
null controls. The four zero 10076 directions require no projections and
remain uninformative. Only one CTF/noise profile is studied. We use
$\alpha=.05$, $\delta=.001$, $\kappa\in\{1,1.01,1.1,2\}$ and
$n\in\{10^3,10^4,10^5\}$.

The table projects dataset-level rejection probabilities from independent
single-particle event frequencies under the binomial model. Its intervals
transform pointwise exact 95\% probability intervals; they are not observed
repeated-dataset power or simultaneous confidence bounds across all comparisons.
The complete CSV retains both candidate controls and every sample size.
At $n=10^4$ and $\kappa=1$, ranged-amplitude 10049 rejection projections
are .700 for power and .826 for combined moments, with overlapping pointwise
intervals [.492,.857] and [.654,.930]. The two ranged-amplitude 10028 tests
instead have projections below .01. At $n=10^5$, all 10049 projections
collapse below $10^{-14}$ when $\kappa=1.1$; increasing $\kappa$ changes the
critical bound while the held-out alternative remains Haar. This measures
conservatism of this construction, not an impossibility result or an observed
non-Haar experiment. Correct-null point projections range up to .0435; they
do not measure unconditional false-rejection frequency over calibration draws.
Scores were designed with the oracle alternative, so this experiment does not
establish detection of unknown structural errors. A known viewing-density
bound, calibrated noise/amplitude law, simulator error and practical candidate
construction are additional unresolved requirements. These results cannot be
substituted for experimental density coverage or the failed arbitrary-view
certificate.
